% ============================================================================
% Fichier  : partie3/P03_C10_cadre_normatif_global.tex
% Titre    : Cadre Normatif Global
% Auteur   : MINKA MI NGUIDJOÏ Thierry Emmanuel
% Version  : 1.0 -- Juin 2026
% Statut   : RÉDIGÉ
% Sources  : 5_cadre_normatif_global.tex (squelette ISO)
%            M4_Cadre_Normatif_International_PNC.docx (Budapest, Malabo,
%            INTERPOL, AFRIPOL, portails LEA, exercices BAOBAB)
%            Ch. 7 (standards forensiques déjà couverts)
% Positionnement : Ch. 7 = standards du point de vue technique-forensique.
%                  Ce chapitre = cadre normatif du point de vue coopératif
%                  et de la légitimité internationale des investigations.
% ============================================================================

\chapter{Cadre Normatif Global}
\label{chap:norm-glob}

\epigraph{%
    « La loi doit fournir un cadre sans entraver l'innovation,
    protéger sans étouffer, réguler sans paralyser. »%
}{--- Simone Veil}

\niveauChapitre{Fondamental}{%
    Aucun prérequis technique. Ce chapitre est essentiellement juridique
    et procédural. Lecture du Chapitre~\ref{chap:meth-std} (Standards)
    recommandée pour ne pas dupliquer --- les deux se complètent.%
}

\begin{boxobjectifs}
\begin{itemize}
    \item Situer la Convention de Budapest dans le paysage du droit
          international du cyber et comprendre sa relation avec la
          législation camerounaise.
    \item Maîtriser les articles fondamentaux de Budapest~: art.~2--11
          (infractions), art.~16--21 (mesures procédurales),
          art.~23--35 (coopération).
    \item Comparer Budapest et Malabo sur cinq critères~: champ d'application,
          infractions, protection des données, coopération, ratifications.
    \item Mobiliser le bon instrument international selon la nature,
          l'urgence et la géographie d'une demande de coopération.
    \item Utiliser les canaux INTERPOL (réseau 24/7) et AFRIPOL
          pour des demandes de coopération urgentes.
    \item Soumettre une demande structurée aux portails LEA des grands
          fournisseurs (Meta, Google, Microsoft).
\end{itemize}
\end{boxobjectifs}

\bigskip

% ============================================================================
\section*{Pourquoi un cadre normatif global~?}
% ============================================================================

Le cybercrime n'a pas de frontières. Un ransomware déployé depuis
Saint-Pétersbourg, chiffrant des fichiers à Yaoundé, via un serveur de
commande en Roumanie, est une réalité quotidienne. Sans cadre normatif
commun entre les États, chaque frontière est un mur~: les preuves
disparaissent avant que les demandes d'entraide n'arrivent, les suspects
se réfugient dans des pays sans traités, les fonds blanchis traversent
des juridictions incompatibles.

Ce chapitre répond à une question pratique~: \emph{quand une affaire
dépasse les frontières camerounaises, sur quelle base légale demander
quoi, à qui, par quel canal, dans quel délai~?}

\separateur

% ============================================================================
\section{La Convention de Budapest~: Premier Traité International
sur la Cybercriminalité}
\label{sec:budapest}
% ============================================================================

\subsection{Contexte, portée et position du Cameroun}

Adoptée par le Conseil de l'Europe le 23~novembre~2001 (STE n°~185),
la \termecle{Convention de Budapest}\index{Budapest, Convention de} est
le premier et unique traité international contraignant en matière de
cybercriminalité. En 2026, plus de 68~États la ont ratifiée, dont les
États-Unis, le Japon, l'Australie, le Canada, et sur le continent
africain~: le Cap-Vert, le Ghana, le Maroc, la Mauritanie, le Nigeria,
le Sénégal et la Tunisie.

\begin{boxalert}
\textbf{Position du Cameroun vis-à-vis de Budapest~:}

Le Cameroun \textbf{n'est pas signataire} de la Convention de Budapest
(juin~2026). Ce fait ne le place pas hors du droit international du
cyber pour autant~:
\begin{enumerate}
    \item La \loicam{2010/012} s'est directement inspirée de Budapest~---
          les deux textes sont largement alignés sur les infractions et
          les mesures procédurales.
    \item Les partenaires africains ayant ratifié Budapest (Ghana, Nigeria,
          Sénégal) coopèrent avec le Cameroun via leurs propres obligations.
    \item Le Cameroun peut rejoindre Budapest à tout moment --- la Convention
          est ouverte à tous les États.
\end{enumerate}
\textbf{Posture opérationnelle~:} maîtriser Budapest comme si le Cameroun
en était partie --- parce que les partenaires avec qui les OPJ coopèrent
l'appliquent, et parce que la Loi~2010/012 en est directement issue.
\end{boxalert}

\subsection{Architecture de la Convention}

\begin{tableaucours}{p{4.0cm}p{8.0cm}}{%
    \tableauHeader{Section} &
    \tableauHeader{Contenu essentiel}
}
Chapitre~I --- Définitions (art.~1) &
    Lexique commun~: système informatique, données informatiques,
    fournisseur de services, données de trafic. \\
Chapitre~II Sect.~1 --- Droit pénal (art.~2--13) &
    Infractions que chaque État doit ériger en droit interne. \\
Chapitre~II Sect.~2 --- Procédure (art.~14--22) &
    Pouvoirs d'investigation~: conservation, production,
    perquisition, interception. \\
Chapitre~III --- Coopération (art.~23--35) &
    Mécanismes d'entraide~: extradition, MLAT, réseau 24/7,
    conservation urgente transfrontalière. \textbf{Le cœur opérationnel.} \\
Chapitre~IV --- Dispositions finales (art.~36--48) &
    Signatures, ratifications, réserves. \\
\end{tableaucours}
\vspace{-0.8em}
\begin{center}{\small\itshape Tableau~10.1 --- Architecture de la Convention de Budapest}\end{center}

% ============================================================================
\subsection{Les infractions pénales~: articles~2 à~11}
% ============================================================================

Ces articles définissent les infractions que chaque État partie doit
ériger en droit interne. Pour un OPJ camerounais, ils fondent les
demandes de coopération vers les États partenaires.

\begin{boxjuris}
\textbf{Article~2 --- Accès illégal~:}
\begin{quote}
Chaque Partie adopte les mesures législatives \textit{[\ldots]} pour
ériger en infraction pénale \textit{[\ldots]} le fait intentionnel
d'accéder sans droit à tout ou partie d'un système informatique.
\end{quote}
\textbf{Équivalent camerounais~:} \loicam{2010/012}, art.~55.
Formulations quasi-identiques.\\
\textbf{Usage opérationnel~:} art.~2 fonde la demande de coopération
quand le serveur piraté ou l'auteur de l'intrusion est situé dans un
État partie à Budapest.
\end{boxjuris}

\begin{boxjuris}
\textbf{Article~8 --- Fraude informatique~:}
\begin{quote}
\textit{[\ldots]} le fait de causer un préjudice patrimonial à autrui
par toute introduction, altération, effacement ou suppression de données
informatiques \textit{[\ldots]} dans l'intention frauduleuse ou déloyale
d'obtenir sans droit un bénéfice économique.
\end{quote}
\textbf{Équivalent camerounais~:} \loicam{2010/012}, art.~78.
\textbf{L'article le plus utilisé en pratique.}\\
Couvre~: BEC, phishing bancaire, fraude Mobile Money, virement
frauduleux par manipulation. Art.~8 fonde la demande vers le pays de
l'auteur \emph{ET} vers la banque destinataire pour gel des fonds.
\end{boxjuris}

\begin{tableaucours}{p{1.8cm}p{3.5cm}p{2.8cm}p{3.2cm}}{%
    \tableauHeader{Article} &
    \tableauHeader{Infraction} &
    \tableauHeader{Équiv. Loi~2010} &
    \tableauHeader{Cas typique Cameroun}
}
Art.~2 & Accès illégal & Art.~55 & Intrusion CNPS, système bancaire \\
Art.~3 & Interception illégale & Art.~56 & Logiciel espion C2 étranger \\
Art.~4 & Atteinte aux données & Art.~63 & Ransomware hôpital \\
Art.~5 & Atteinte au système & Art.~65 & DDoS infrastructure \\
Art.~6 & Abus de dispositifs & Art.~67 & Kit phishing vendu sur dark web \\
Art.~7 & Falsification & Art.~71 & Documents administratifs falsifiés \\
Art.~8 & Fraude informatique & Art.~78 & BEC, MoMo fraud, phishing \\
Art.~9 & CSAM & Art.~76 & Contenus pédopornographiques \\
\end{tableaucours}
\vspace{-0.8em}
\begin{center}{\small\itshape Tableau~10.2 --- Articles Budapest, équivalents camerounais
et cas typiques}\end{center}

% ============================================================================
\subsection{Les mesures procédurales~: articles~16, 17, 18 et 29}
% ============================================================================

Ces articles sont les plus directement utiles à l'enquêteur. Ils
définissent les outils d'investigation que chaque État doit mettre
à disposition et les mécanismes de coopération procédurale.

\begin{boxjuris}
\textbf{Article~16 --- Conservation rapide de données stockées~:}
\begin{quote}
\textit{[\ldots]} ordonner ou d'imposer \textit{[\ldots]} la conservation
rapide de données électroniques spécifiées \textit{[\ldots]} notamment
lorsqu'il y a des raisons de croire que celles-ci sont particulièrement
susceptibles de perte ou de modification. La période de conservation est
aussi longue que nécessaire, jusqu'à un maximum de 90~jours.
\end{quote}
\textbf{L'outil d'urgence n°~1 de l'enquêteur cyber.} Les données
numériques disparaissent vite~: logs effacés, comptes supprimés, serveurs
réinitialisés.\\[0.3ex]
\textbf{Conservation $\neq$ Divulgation~:} art.~16 préserve les données
(mise sous scellé numérique temporaire). La divulgation nécessite
art.~17--18. Ces deux étapes ne peuvent pas être sautées.
\end{boxjuris}

\begin{boxjuris}
\textbf{Article~29 --- Conservation urgente transfrontalière~:}
\begin{quote}
Une Partie peut demander à une autre Partie d'ordonner ou d'imposer
d'une autre façon la conservation rapide de données stockées au moyen
d'un système informatique situé sur le territoire de cette autre
Partie et pour lesquelles la Partie requérante souhaite présenter une
demande d'entraide en vue d'obtenir leur divulgation.
\end{quote}
\textbf{Réflexe obligatoire~:} dès qu'une adresse IP suspecte est
identifiée sur un serveur étranger, envoyer \emph{immédiatement} une
demande de conservation via le réseau 24/7 INTERPOL. Attendre la voie
diplomatique classique = arriver après effacement des données.
\end{boxjuris}

\begin{boxjuris}
\textbf{Article~35 --- Réseau 24/7~:}
Chaque Partie désigne un point de contact disponible 24h/24, 7j/7,
pour assister les autres Parties dans les investigations sur les
systèmes informatiques et les données stockées.\\[0.3ex]
Pour le Cameroun~: \textbf{BCN INTERPOL Yaoundé (DGSN/DCPJ)} est le
point de contact pour les demandes entrantes et sortantes via le réseau
I-24/7.
\end{boxjuris}

\begin{boxCRO}
\textbf{Analyse CRO --- Article~29 et conservation urgente}

$\mathcal{R}$ (Fiabilité)~: la conservation urgente préserve les données
avant qu'elles ne disparaissent --- elle maximise la fiabilité future
du dossier.

$\mathcal{O}$ (Opposabilité)~: art.~29 fonde légalement la demande~;
sans cette base, les données obtenues peuvent être contestées comme
ayant été recueillies sans fondement.

$\mathcal{C}$ (Confidentialité)~: la conservation implique que le
fournisseur étranger a connaissance de l'investigation. Dans certains
cas, cela peut alerter le suspect (via compromission du fournisseur).

\cro{$\downarrow$ risque à évaluer}{$\uparrow\uparrow$ conservation = intégrité future}{$\uparrow\uparrow$ base légale solide}
\end{boxCRO}

% ============================================================================
\section{La Convention de Malabo~: Réponse Africaine}
\label{sec:malabo}
% ============================================================================

\subsection{Contexte et entrée en vigueur}

La \termecle{Convention de Malabo}\index{Malabo, Convention de} (Convention
de l'Union Africaine sur la cybersécurité et la protection des données
personnelles) a été adoptée le 27~juin~2014 et est entrée en vigueur
le 8~juin~2023, après sa 15\textsuperscript{e} ratification.
En~2026, environ 15--16~États africains l'ont ratifiée.

Contrairement à Budapest (purement pénale), Malabo est une convention
\textbf{trifonctionnelle}~: elle couvre simultanément le commerce
électronique (Chap.~II), la protection des données personnelles
(Chap.~III) et la cybercriminalité (Chap.~V).

\subsection{Infractions spécifiques à Malabo}

Par rapport à Budapest, Malabo ajoute trois infractions absentes~:

\begin{itemize}
    \item \textbf{Usurpation d'identité (art.~V-41)}~: ABSENTE de Budapest.
          Malabo criminalise spécifiquement l'usurpation d'identité en
          ligne. Directement applicable~: SIM~swapping, faux profils
          Facebook, phishing par usurpation de marque.
    \item \textbf{Cyberterrorisme (art.~V-42)}~: actes terroristes via
          les TIC ou ciblant des infrastructures critiques. Pertinent
          pour les menaces Boko~Haram utilisant les réseaux sociaux.
    \item \textbf{Cybersquatting (art.~V-46)}~: enregistrement de noms
          de domaine imitant des marques ou institutions de mauvaise foi~;
          faux sites imitant les banques camerounaises, administrations,
          opérateurs Mobile Money.
\end{itemize}

\subsection{Tableau comparatif Budapest / Malabo}

\begin{tableaucours}{p{3.5cm}p{4.5cm}p{4.5cm}}{%
    \tableauHeader{Critère} &
    \tableauHeader{Budapest (2001)} &
    \tableauHeader{Malabo (2014)}
}
Origine & Conseil de l'Europe & Union Africaine \\
Portée & Mondiale (ouverte à tous) & Africaine (membres UA) \\
États parties (2026) & $\approx$~68 & $\approx$~15--16 \\
Entrée en vigueur & 1\textsuperscript{er} juillet~2004 & 8~juin~2023 \\
Infractions spécif. & Art.~2--11 + CSAM + droit d'auteur &
    Idem + usurpation identité + cyberterrorisme + cybersquatting \\
Protection données & Absente (couverte par Convention~108) &
    Intégrée (Chap.~III complet) \\
Commerce électronique & Absent & Intégré (Chap.~II) \\
Coopération & Réseau 24/7 (art.~35), MLAT, art.~29 & AFRIPOL~; moins formalisé \\
Force contraignante & Élevée~; jurisprudence abondante &
    Plus faible~; texte récent, peu de pratique \\
Position Cameroun & Non signataire --- influence via Loi~2010 &
    Signataire non-ratifiant (juin~2026) \\
\end{tableaucours}
\vspace{-0.8em}
\begin{center}{\small\itshape Tableau~10.3 --- Comparaison Budapest / Malabo}\end{center}

\begin{boiteRemarque}{Règles de choix entre Budapest et Malabo}
\begin{enumerate}
    \item Utiliser toujours le texte que l'État requis a ratifié.
    \item Si l'État requis a ratifié les deux~$\to$ Budapest en priorité
          (plus de pratique, mécanismes plus rodés).
    \item Si l'État requis a ratifié Malabo mais pas Budapest
          (ex.~: Angola)~$\to$ Malabo Chap.~V.
    \item Si l'État requis n'a ratifié ni l'un ni l'autre~$\to$
          accords bilatéraux ou canal INTERPOL/AFRIPOL.
    \item Vérification obligatoire avant toute demande~:
          \texttt{treaty.un.org} (Budapest) et \texttt{au.int} (Malabo).
\end{enumerate}
\end{boiteRemarque}

% ============================================================================
\section{Les Protocoles Additionnels de Budapest}
\label{sec:protocoles-budapest}
% ============================================================================

\begin{tableaucours}{p{3.5cm}p{4.0cm}p{4.5cm}}{%
    \tableauHeader{Protocole} &
    \tableauHeader{Objet} &
    \tableauHeader{Impact opérationnel}
}
\textbf{Protocole~I} (STE~189, 2003) &
    Contenus racistes et xénophobes &
    35~États. USA non-signataire.
    Demandes difficiles vers les USA pour ce type de contenu. \\
\textbf{Protocole~II} (STE~224, 2022,
    en vigueur~2024) &
    Accès élargi aux données stockées dans le cloud &
    Délai 45--90~jours vs 18~mois pour MLAT classique.
    Révolutionne la forensique cloud transfrontalière. \\
\end{tableaucours}
\vspace{-0.8em}
\begin{center}{\small\itshape Tableau~10.4 --- Protocoles additionnels Budapest}\end{center}

Le Protocole~II est particulièrement important~: il crée une procédure
accélérée pour accéder aux données cloud détenues par des fournisseurs
dans les États parties. Pour la forensique cloud, il représente un
changement de paradigme~: délais de mois réduits à semaines.

% ============================================================================
\section{INTERPOL et AFRIPOL~: Architecture de la Coopération}
\label{sec:interpol-afripol}
% ============================================================================

\subsection{INTERPOL --- réseau mondial}

\begin{tableaucours}{p{3.5cm}p{8.0cm}}{%
    \tableauHeader{Composante} &
    \tableauHeader{Description}
}
Structure & 195~pays membres. Siège~: Lyon. Bureau régional Afrique
    centrale~: Yaoundé. Direction Cybercriminalité~: Singapour (IGCI). \\
Réseau I-24/7 & Réseau mondial sécurisé entre les 195~BCN.
    Accessible uniquement via le BCN national
    (DGSN pour le Cameroun). \\
Notices & 7~types~: \textbf{Rouge} (arrestation), \textbf{Bleue}
    (renseignements), Verte (avertissement), Mauve (modes opératoires),
    Orange (menaces), Jaune (disparitions), Noire (cadavres). \\
Opérations récentes & HAECHI~I--IV (fraudes en ligne Asie)~;
    FALCON~I--II (cybercriminels africains, Cameroun participant)~;
    GOLDFISH~ALPHA (malwares routeurs). \\
\end{tableaucours}
\vspace{-0.8em}
\begin{center}{\small\itshape Tableau~10.5 --- Architecture INTERPOL cyber}\end{center}

\subsubsection*{Procédure d'activation du réseau 24/7}

\begin{tableaucours}{p{3.5cm}p{8.0cm}}{%
    \tableauHeader{Paramètre} &
    \tableauHeader{Détail}
}
Qui peut activer~? &
    Tout OPJ, via la chaîne hiérarchique~$\to$ BCN Cameroun
    (DGSN/DCPJ)~$\to$ réseau I-24/7. Impossible de bypasser le BCN. \\
Types de demandes &
    Conservation urgente (art.~29 Budapest), identification IP,
    identification abonné, localisation suspect. \\
Délais moyens &
    Conservation~: 24--48h. Identification IP~: 48--72h.
    Identification abonné~: 72h--7~jours. \\
Format &
    Fiche INTERPOL standardisée (disponible au BCN). Minimum~:
    infraction, données à identifier, urgence motivée, coordonnées
    de contact. \\
Suivi &
    Toujours noter un numéro de référence interne. Relancer après
    48h sans réponse. Urgence vitale~: appel téléphonique direct
    au BCN du pays requis. \\
\end{tableaucours}
\vspace{-0.8em}
\begin{center}{\small\itshape Tableau~10.6 --- Procédure réseau 24/7 INTERPOL}\end{center}

\subsection{AFRIPOL --- coopération africaine}

\termecle{AFRIPOL}\index{AFRIPOL} est le mécanisme de coopération policière
de l'Union Africaine~: opérationnel depuis 2014, siège à Alger, unité
Cybercriminalité créée en~2017.

\begin{itemize}
    \item \textbf{Opération FALCON~I (2021)}~: 11~pays africains~;
          11~cybercriminels arrêtés (groupe SilverTerrier, Nigeria)~;
          50~000~appareils infectés identifiés.
    \item \textbf{Opération FALCON~II (2022)}~: 14~pays africains
          \textbf{dont le Cameroun}~; ciblage des réseaux BEC et
          phishing transfrontaliers.
    \item \textbf{Avantage comparatif}~: meilleure compréhension du
          contexte africain, moins de barrières culturelles. Utiliser
          en complément d'INTERPOL, pas à sa place.
\end{itemize}

% ============================================================================
\section{Les MLAT et les Portails LEA}
\label{sec:mlat-lea}
% ============================================================================

\subsection{La règle d'or~: conservation avant divulgation}

\begin{boxalert}
\textbf{RÈGLE D'OR --- CONSERVATION AVANT TOUT}
\begin{enumerate}
    \item \textbf{Conservation IMMÉDIATE} --- réseau 24/7 INTERPOL /
          portail LEA urgence --- délai cible~: 24--72h.
    \item \textbf{Qualification + instrument} --- Budapest / Malabo /
          accord bilatéral applicable.
    \item \textbf{MLAT formel} --- dans les 60~jours --- pour obtenir
          les données avec valeur probatoire opposable en tribunal.
\end{enumerate}
\textbf{Erreur fatale~:} attendre le MLAT pour demander la conservation.
Le MLAT prend 3~semaines à 18~mois. Les données volatiles disparaissent
en heures.
\end{boxalert}

\subsection{Les portails LEA des grands fournisseurs}

\begin{tableaucours}{p{3.0cm}p{3.5cm}p{3.0cm}p{2.5cm}}{%
    \tableauHeader{Fournisseur} &
    \tableauHeader{Portail} &
    \tableauHeader{Données disponibles} &
    \tableauHeader{Délais}
}
\textbf{Meta} (FB, IG, WA) &
    \texttt{facebook.com/records/login} &
    Abonné, logs, IP connexion. Contenu messages~: MLAT USA. &
    Standard~: 7--21j. Urgence~: 24--72h. \\
\textbf{Google} (Gmail, YT, Android) &
    \texttt{support.google.com/legal} &
    Abonné Gmail, localisation Android (précise). Contenu~: MLAT. &
    Standard~: 10--30j. Urgence~: 24--48h. \\
\textbf{Microsoft} (Outlook, OneDrive) &
    \texttt{microsoft.com/legal/lerp} &
    Abonné, historique connexion. Contenu~: MLAT. &
    Le plus réactif aux demandes étrangères. \\
\textbf{TikTok} &
    \texttt{tiktok.com/legal/law-enforcement} &
    Abonné, activité. &
    Standard~: 14--30j. \\
\textbf{Telegram} &
    (portail limité) &
    Répond rarement sans ordonnance. &
    Privilégier MLAT. \\
\end{tableaucours}
\vspace{-0.8em}
\begin{center}{\small\itshape Tableau~10.7 --- Portails LEA des grands fournisseurs}\end{center}

\begin{boxPNC}
\textbf{WhatsApp~: limite technique irréductible}

Le chiffrement de bout en bout d'WhatsApp signifie que Meta \emph{ne peut
pas} fournir le contenu des messages, même avec une ordonnance. Accessible
via portail LEA~: uniquement les métadonnées (qui a communiqué avec qui,
quand, durée). Le contenu des messages nécessite un MLAT USA --- et même
alors, Meta ne peut fournir que ce qu'elle stocke (métadonnées).

\textbf{Alternative pratique~:} saisir physiquement le téléphone du suspect
et acquérir les messages via ADB (si téléphone déverrouillé) --- la donnée
est sur l'appareil, pas chez Meta.
\end{boxPNC}

% ============================================================================
\section{Tableau de Décision~: Quel Canal pour Quelle Demande~?}
\label{sec:tableau-decision}
% ============================================================================

\begin{tableaucours}{p{3.5cm}p{4.5cm}p{4.0cm}}{%
    \tableauHeader{Situation} &
    \tableauHeader{Protocole d'action} &
    \tableauHeader{Délais et remarques}
}
Urgence vitale (menace de mort, CSAM actif) &
    1.~Portail LEA urgent (24--72h). 2.~BCN INTERPOL 24/7.
    3.~Conservation urgente art.~29. &
    Agir immédiatement. Ne pas attendre l'autorisation
    judiciaire formelle pour la conservation. \\
Données volatiles (IP suspecte, logs) &
    1.~\emph{IMMÉDIATEMENT}~: conservation BCN INTERPOL.
    2.~Portail LEA en parallèle. 3.~MLAT dans les 60~jours. &
    Conservation = réflexe~n°~1. Tout le reste peut attendre. \\
Identification IP étrangère &
    1.~WHOIS (pays propriétaire). 2.~Si Budapest~: réseau 24/7 BCN.
    3.~Réquisition FAI camerounais si IP de transit. &
    48h--7j. FAI européens/américains~: conservation 6~mois--1~an. \\
Compte réseau social &
    1.~Portail LEA fournisseur. 2.~Urgence~: cocher Emergency.
    3.~BCN INTERPOL si délai insuffisant. &
    Abonné~: 7--30j. Urgence~: 24--72h. Contenu~: MLAT (mois). \\
Gel de fonds crypto &
    1.~Identifier plateforme et pays. 2.~GABAC~$\to$ GIABA.
    3.~Demande gel via canaux financiers. &
    Plateformes légales~: obligations AML/KYC.
    Bitcoin~: analyse blockchain via Blockchair. \\
État africain signataire Malabo &
    1.~Vérifier ratification (au.int). 2.~Citer Malabo Chap.~V.
    3.~AFRIPOL si nécessaire. &
    En cas de doute~: contacter AFRIPOL directement. \\
État non-signataire Budapest/Malabo &
    1.~Accords bilatéraux. 2.~Canal INTERPOL hors Budapest.
    3.~Chercher nœuds de la chaîne dans des pays coopérants. &
    Russie, Chine~: coopération limitée. Chercher les nœuds
    intermédiaires dans des pays plus accessibles. \\
\end{tableaucours}
\vspace{-0.8em}
\begin{center}{\small\itshape Tableau~10.8 --- Tableau de décision~: canal selon la situation}\end{center}

% ============================================================================
\section{Cas Pratique~: Opération BAOBAB}
\label{sec:baobab}
% ============================================================================

\begin{boxscenario}{Opération BAOBAB --- Fraudes MoMo transfrontalières}
\textbf{Contexte~:} Fraudes MTN MoMo~: 47~millions FCFA détournés
en 3~mois sur 23~victimes camerounaises. Les fonds ont été transférés
vers un compte Binance, convertis en Bitcoin et envoyés vers une adresse
dont la dernière connexion provient d'une IP sénégalaise
(\texttt{41.82.XXX.XXX} --- Sonatel).

\textbf{Données disponibles~:}
\begin{enumerate}
    \item IP de connexion au compte Binance.
    \item Email du compte Binance.
    \item Adresse Bitcoin de destination.
    \item Dates et montants de toutes les transactions MTN MoMo.
\end{enumerate}

\textbf{Trois actions simultanées~:}
\end{boxscenario}

\begin{tableaucours}{p{3.2cm}p{5.0cm}p{4.0cm}}{%
    \tableauHeader{Action} &
    \tableauHeader{Description} &
    \tableauHeader{Base légale et délai}
}
\textbf{Action~1} --- Conservation Binance &
    Demande de conservation urgente~: données d'abonne, logs,
    transactions. Via portail LEA Binance + BCN INTERPOL. &
    Art.~16/29 Budapest si État partie~; portail LEA. Cible~: 48h. \\
\textbf{Action~2} --- Identification IP Sénégal &
    Demande d'entraide urgente vers Sonatel pour identification
    de l'abonné de l'IP sénégalaise. &
    Art.~29 Budapest (Sénégal est partie). BCN Sénégal. 7--14~jours. \\
\textbf{Action~3} --- Analyse blockchain &
    Tracer l'adresse Bitcoin de destination. Identifier les transferts
    ultérieurs et les plateformes d'échange. &
    \texttt{Blockchair.com} (gratuit, public). Résultat immédiat. \\
\end{tableaucours}
\vspace{-0.8em}
\begin{center}{\small\itshape Tableau~10.9 --- Opération BAOBAB~: trois actions simultanées}\end{center}

\begin{boxalert}
\textbf{Erreurs les plus fréquentes dans cet exercice~:}
\begin{enumerate}
    \item Envoyer directement une demande de divulgation sans conservation
          préalable~: les données sont effacées avant réponse.
    \item Ne pas citer la base légale dans la demande~: une demande sans
          base légale est systématiquement renvoyée.
    \item Adresser la demande directement à Sonatel sans canal officiel~:
          une demande directe à un opérateur étranger sans voie officielle
          a peu de chances d'aboutir rapidement.
\end{enumerate}
\end{boxalert}

\separateur

\begin{boxresume}
\textbf{Ce qu'il faut retenir de ce chapitre~:}
\begin{itemize}
    \item Le Cameroun \textbf{n'est pas signataire de Budapest} mais la
          Loi~2010/012 en est directement issue~: les deux textes sont
          alignés. Maîtriser Budapest comme si le Cameroun en était partie.

    \item \textbf{Articles Budapest à retenir~:} art.~2 (accès illégal)~;
          art.~8 (fraude~: le plus utilisé)~; art.~16 (conservation
          nationale~: 90~jours max)~; art.~29 (conservation urgente
          transfrontalière~: réflexe~n°~1)~; art.~35 (réseau 24/7).

    \item \textbf{Malabo} ajoute trois infractions absentes de Budapest~:
          usurpation d'identité, cyberterrorisme, cybersquatting. Moins
          de pratique mais adapté à la réalité africaine.

    \item \textbf{Règle d'or}~: Conservation~$\to$ Qualification~$\to$
          MLAT (dans cet ordre, toujours). Ne jamais attendre le MLAT
          pour demander la conservation.

    \item \textbf{Portails LEA}~: Meta/Google/Microsoft permettent
          d'obtenir les données d'abonné en 24h--30j. Le contenu des
          messages nécessite toujours un MLAT.

    \item \textbf{INTERPOL réseau 24/7}~: accessible via le BCN
          Cameroun (DGSN). Délais~: 24--48h pour conservation,
          48--72h pour identification IP.

    \item \textbf{Le tableau de décision} (Tableau~10.8) est le mémo
          opérationnel~: une ligne par situation, trois colonnes~:
          protocole, canal, délai.
\end{itemize}
\end{boxresume}

\bigskip

\begin{boxexo}[title={\ding{45}\quad Exercice~10.1 --- Qualification internationale \niveaubase}]
Pour chacune des situations suivantes, identifiez~: (A)~l'article Budapest
applicable, (B)~l'article équivalent de la Loi~2010/012, (C)~une divergence
éventuelle et ses conséquences pratiques~:
\begin{enumerate}[(a)]
    \item Un groupe de hackers depuis le Nigeria a accédé au système de
          gestion des marchés publics camerounais et modifié des données
          d'attribution.
    \item Un ressortissant sénégalais a intercepté les communications
          d'un dirigeant d'entreprise camerounais via un logiciel espion.
    \item Un réseau distribué dans 12~pays envoie 4~millions de faux
          emails au nom d'UBA Cameroun, récoltant les identifiants
          bancaires de 3~000~victimes.
    \item Un site web hébergé en Côte~d'Ivoire imite fidèlement le portail
          de Mobile Money de MTN Cameroun pour voler les identifiants.
\end{enumerate}
\end{boxexo}

\begin{boxexo}[title={\ding{45}\quad Exercice~10.2 --- Choix de canal \niveaumoyen}]
Pour chacun des 6~scénarios suivants, identifiez le canal de coopération
à activer en premier, la base légale applicable, et le délai attendu~:
\begin{enumerate}[(a)]
    \item Serveur de phishing bancaire hébergé aux Pays-Bas (Budapest).
    \item Compte Facebook d'un harceleur identifié au Sénégal
          (Budapest + Malabo).
    \item Fonds blanchis sur Orange Money Côte~d'Ivoire (Malabo ratifié).
    \item Serveur C2 de logiciel espion situé en Russie
          (non-Budapest, non-Malabo).
    \item Vidéo terroriste Boko~Haram sur YouTube (serveurs USA).
    \item Contenu CSAM sur Telegram.
\end{enumerate}
\end{boxexo}

\begin{boxexo}[title={\ding{45}\quad Exercice~10.3 --- Rédaction de demande de conservation \niveauexpert}]
Dans l'Opération WOURI (Chapitre~\ref{chap:gr-affaires}), vous avez
identifié~: l'IP \texttt{91.XXX.XX.XX} (enregistrée aux Pays-Bas) comme
origine de l'email de phishing~; un compte Binance utilisé pour le
blanchiment~; une adresse Bitcoin liée à des transactions vers un exchange
aux Émirats arabes unis (Budapest-ratifié depuis 2021).

\begin{enumerate}[(a)]
    \item Rédigez la demande de conservation urgente selon art.~29 Budapest
          pour l'IP néerlandaise~: base légale, nature de l'infraction,
          données demandées, urgence motivée, délai demandé.
    \item Quels sont les trois éléments \emph{sans lesquels} la demande
          sera systématiquement rejetée~?
    \item Appliquez la grille CRO à la décision d'activer simultanément
          les trois canaux (Pays-Bas, Binance, EAU)~: quels risques
          pour $\mathcal{C}$ et comment les gérer~?
\end{enumerate}
\end{boxexo}

\bigskip

\begin{boxinfo}
\textbf{Transition.} Le Chapitre~\ref{chap:leg-monde} approfondit la
\textbf{législation mondiale comparative}~: RGPD, CLOUD~Act, NIS~2 en
Europe~; CCPA aux États-Unis~; PIPL en Chine. Ces législations
conditionnent directement la coopération internationale~: savoir ce que
le droit local d'un État partenaire autorise ou interdit est
indispensable pour formuler des demandes qui aboutissent.
\end{boxinfo}
